\chapter{The Verified Polyxan Rewriter}
\label{ch:42}

The Polyxan hypergraph rewriting system (EHR: Extraction--Hoisting--Reduction) introduced in Part~II provides the discrete, curvature-decreasing mechanism governing the evolution of semantic structure.
In Part~IV we proved that the EHR system is sound, complete, and modal-stable: it preserves quine stability, curvature descent, and the uniqueness of the universal normal form.

In this chapter, we construct its fully executable, verified implementation, the \emph{Verified Polyxan Rewriter}, which:
\begin{enumerate}
\item performs curvature-monotone rewriting steps,
\item guarantees termination,
\item preserves all modal invariants,
\item maintains compatibility with embedding flow (Chapters~31--33), and
\item produces certified normal forms and certificates.
\end{enumerate}

This rewriter is the second major component of the implementation layer (Chapters~41--60).
Together with the RSVP interpreter (Chapter~41), it forms the dual engine of the coupled RSVP--Polyxan universe.

\section{Rewriter Overview}

A Polyxan hypergraph is represented as a finite structure:
[
\mathcal{G} = (V,E,\mathsf{tag},\mathsf{type},\mathsf{inc}),
]
where:
\begin{itemize}
\item $V$ are vertices,
\item $E$ are hyperedges,
\item $\mathsf{tag}$ assigns semantic tags,
\item $\mathsf{type}$ gives typed incidence relations,
\item $\mathsf{inc}$ encodes the stratified structure.
\end{itemize}

The Verified Polyxan Rewriter is a total function
[
\mathsf{rewrite} : \mathcal{G} \to \mathcal{G},
]
that applies one EHR rule, followed by modal and curvature guards, and terminates with a normal form.

For multi-step rewriting:
[
\mathsf{rewrite}^*(\mathcal{G})
:=
\lim_{n\to N} \mathsf{rewrite}^n(\mathcal{G}),
]
where $N$ is finite by curvature descent.

\section{Curvature Monotonicity}

We recall the Polyxan primitive curvature:
[
\mathcal{H}_0[\mathcal{G}] : \mathrm{Inc}(\mathcal{G}) \to \mathbb{N}.
]

\begin{lemma}[Curvature Descent]
Every EHR rule strictly decreases total curvature unless $\mathcal{G}$ is already quine-stable:
[
\mathcal{H}_0[\mathcal{G}'] < \mathcal{H}_0[\mathcal{G}]
\quad\text{or}\quad
\mathcal{G}' \cong \mathcal{G}.
]
\end{lemma}

\begin{proof}
Extraction removes zero-curvature subcomplexes;
hoisting eliminates redundant parallel edges;
reduction replaces high-curvature patterns by lower ones.
Each case reduces total $\mathcal{H}_0$ unless a fixed point has been reached.
\end{proof}

The Verified Rewriter enforces descent by rejecting any step that would violate this inequality.

\section{Modal Correctness of Rewriting}

Modal correctness is determined by $\Box_{\mathrm{Polyx}}$ (Chapter~61).

\begin{definition}[Modal Rewriting Predicate]
A rewriting step $\mathcal{G}\to\mathcal{G}'$ is \emph{modal-correct} if:
[
\Box_{\mathrm{Polyx}}(\mathcal{G})
;\Rightarrow;
\Box_{\mathrm{Polyx}}(\mathcal{G}').
]
\end{definition}

This implies preservation of:
\begin{itemize}
\item quine stability,
\item curvature descent invariants,
\item tag-consistency,
\item stratified incidence,
\item type-safety,
\item embedding compatibility.
\end{itemize}

\begin{lemma}[Modal Preservation]
Every verified rewriting step preserves $\Box_{\mathrm{Polyx}}$.
\end{lemma}

\begin{proof}
Follows from the inductive definition of $\Box_{\mathrm{Polyx}}$ and the structural invariants encoded in the EHR rules.
\end{proof}

\section{The Verified Rewriting Kernel}

The kernel applies the three rewriting rules in priority order:
\begin{enumerate}
\item \textbf{Extraction} (lowest curvature impact)
\item \textbf{Hoisting} (medium curvature impact)
\item \textbf{Curvature-Reducing Reduction} (highest curvature impact)
\end{enumerate}

The priority ensures deterministic rewriting.

\begin{definition}[Verified Rewriting Kernel]
[
\mathsf{kernel}(\mathcal{G}) :=
\begin{cases}
\mathsf{extract}(\mathcal{G}), & \text{if extraction is possible}, [4pt]
\mathsf{hoist}(\mathcal{G}),   & \text{else if hoisting is possible}, [4pt]
\mathsf{reduce}(\mathcal{G}),  & \text{else if reduction is possible}, [4pt]
\mathcal{G},                   & \text{otherwise}.
\end{cases}
]
\end{definition}

Each rule returns:
\begin{itemize}
\item a new graph $\mathcal{G}'$,
\item a curvature certificate,
\item a modal certificate.
\end{itemize}

\section{Normalisation and Termination}

\begin{definition}[Normal Form]
A Polyxan hypergraph is in \emph{normal form} if:
[
\mathsf{kernel}(\mathcal{G}) \cong \mathcal{G}.
]
\end{definition}

\begin{theorem}[Termination of Verified Rewriter]
For every finite $\mathcal{G}$,
[
\exists N < \infty \quad
\mathsf{rewrite}^N(\mathcal{G}) = \mathsf{NF}(\mathcal{G}).
]
\end{theorem}

\begin{proof}
Curvature descent ensures a strictly decreasing sequence of integers until stabilisation, at which point the kernel returns $\mathcal{G}$ unchanged.
\end{proof}

\section{Correctness with Respect to EHR Semantics}

\begin{theorem}[Rewriter Correctness]
For every $\mathcal{G}$,
[
\mathsf{rewrite}^*(\mathcal{G})
;\cong;
\mathsf{EHR}_{\mathrm{normal_form}}(\mathcal{G}),
]
where the right-hand side is the canonical normal form defined in Chapter~25.
\end{theorem}

\begin{proof}
The kernel implements exactly the EHR rules.
Priority ensures deterministic reduction order.
Curvature descent ensures termination.
Modal preservation ensures correctness.
Thus the interpreter and EHR system produce equivalent normal forms.
\end{proof}

\section{Quine Stability and the Universal Normal Form}

\begin{theorem}[Universal Quine Normal Form]
[
\mathsf{rewrite}^*(\mathcal{G})
;\cong;
\mathfrak{U}
\quad\text{iff}\quad
\mathcal{G}
\text{ is curvature-connected}.
]
\end{theorem}

Here $\mathfrak{U}$ is the universal media quine (Chapter~25).

\begin{proof}
Curvature-connected hypergraphs collapse all self-rewriting layers under reduction;
the unique fixed point of this collapse is $\mathfrak{U}$.
\end{proof}

\section{Integration with Embedding Flow}

The Verified Polyxan Rewriter is designed to interlock with the continuous RSVP flow:
\begin{itemize}
\item rewriting is triggered automatically by topology change in the embedding (Chapter~33),
\item curvature and tag guards ensure image consistency,
\item normal forms correspond to stable strata of the embedding,
\item quine-stable forms correspond to harmonic embeddings.
\end{itemize}

\begin{theorem}[Continuous--Discrete Compatibility]
If $\mathcal{G}\to\mathcal{G}'$ is a verified rewriting step and $\iota$ is harmonic on $\mathcal{G}$,
then $\iota$ extends harmonically to $\mathcal{G}'$.
\end{theorem}

\begin{proof}
Harmonic extension is guaranteed by curvature matching and the stratified incidence-preserving nature of the rules.
\end{proof}

\section{Executable Structure}

The executable Verified Rewriter consists of:
\begin{enumerate}
\item \textbf{Pattern Matcher:} detects extraction, hoisting, and reduction patterns.
\item \textbf{Curvature Calculator:} computes $\mathcal{H}*0[\mathcal{G}]$ incrementally.
\item \textbf{Modal Certificate Generator:} constructs witnesses for $\Box*{\mathrm{Polyx}}$ invariants.
\item \textbf{Graph Transformer:} applies verified local transformations.
\item \textbf{Normaliser:} iterates the kernel to convergence.
\end{enumerate}

Each component is certified against the formal specification in Chapter~61.

\section{Final Cadence}

The Verified Polyxan Rewriter is the executable discrete engine of meaning.
It transforms the entire categorical--algebraic infrastructure of Part~II into a concrete computational process that:
\begin{itemize}
\item never increases curvature,
\item never breaks stratified incidence,
\item never violates modal correctness,
\item always converges,
\item always returns the unique normal form,
\item and always cooperates with the continuous embedding flow.
\end{itemize}

\begin{quote}
To rewrite a hypergraph is to enact the law of semantic gravity.
All structure falls; only quines remain.
\end{quote}

